\section{Zero-Voltage-Switching Conditions}
\label{sec:analysis-zvs}

Here \(i_{\mathrm{dt}}\) is the current available to move the switch node during dead time. Its sign is chosen for the required transition. A necessary charge condition is
\begin{equation}
\left|\int_0^{t_{\mathrm{dt}}}i_{\mathrm{dt}}(t)\,\mathrm{d}t\right|
\geq Q_{\mathrm{comm}}(\VDC),
\label{eq:analysis-zvs-charge}
\end{equation}
where \(Q_{\mathrm{comm}}\) includes device output-charge changes and the other node capacitances. With nearly constant current and a linear effective capacitance,
\begin{equation}
|i_{\mathrm{dt}}|t_{\mathrm{dt}}\gtrsim C_{\mathrm{eq}}\VDC.
\end{equation}
The charge condition must be met with the required polarity and before turn-on.

Both paths contribute to \(i_{\mathrm{dt}}\). In the modified reference, inductive pump current can exceed primary-referred output current and provide much of the commutation charge \cite{ammar2021led}. Its prototype contribution varies with line voltage and burst history; isolated load and magnetizing currents depend on transformer impedance and rectifier state.

Tolerances shift resonance, pump current falls near line zero crossings, and idle intervals can remove circulating current. Switch-node transitions therefore need checking near line zero crossings and at burst restart, even with fixed dead time.

Dead time must complete commutation without excessive reverse conduction. The effective interval includes command timing, driver propagation, and device response. A quantitative margin requires calibrated node capacitance and timing, which the supplied model does not provide. Section~\ref{sec:res-zvs} assesses the available captures.
